\subsection{The adjacent congruence factor: a final bounded source check}

This note checks whether the polarization or the integral congruence
quotient in Berger--Klosin, arXiv:2611.19007v1, supplies the additional
adjacent length required by (D2.14) when $d=2$.  The answer supplied by
the stated results is negative: their inequalities count one adjacent
group and one long-root group.  This is a precise failure of the
attempted application, not an impossibility theorem for a stronger
comparison.  No remaining prime, character, or source point is removed.

\paragraph{The two displayed adjacent ideals are the same ideal.}
In Assumption~2.7 and Lemma~2.9 (printed pp.~6--7), the involution
interchanges the outer blocks and fixes the middle block.  The complete
lemma and Remark~2.10 give
\[
 A_{12}\simeq A_{23},\qquad A_{21}\simeq A_{32},\qquad
 J_{\rm adj}:=T(A_{12}A_{21})=T(A_{23}A_{32}).                 \tag{AF.1}
\]
Consequently the total reducibility ideal is
\[
 I=J_{\rm adj}+T(A_{13}A_{31})+J_{\rm adj}
   =J_{\rm adj}+T(A_{13}A_{31}).                              \tag{AF.2}
\]
The addition here is ideal addition.  Repeating the ideal does not
square it, produce another congruence quotient, or double a length.
The multiplication identities in Remark~2.10 transport this same ideal.
Theorem~10.2 (p.~54) likewise uses complex conjugation to identify the
two adjacent Greenberg groups; it does not exhibit their independent
occurrence in a congruence module.

Nor does this anti-involution make the adjacent product alternating.
For an explicit algebraic check let $\mathcal O$ be a discrete valuation
ring with uniformizer $\pi$, and take the matrix order
\[
 \mathcal E=
 \begin{pmatrix}
 \mathcal O&\pi^2\mathcal O&\pi\mathcal O\\
 \pi\mathcal O&\mathcal O&\pi^2\mathcal O\\
 \pi^2\mathcal O&\pi\mathcal O&\mathcal O
 \end{pmatrix},\qquad \tau(X)=JX^tJ,\quad J_{i,4-i}=1.        \tag{AF.3}
\]
The entrywise valuation inequalities show that this is an algebra;
$\tau$ preserves it and satisfies Assumption~2.7.  Here
$A_{12}=A'_{12}$, $A_{13}/A'_{13}=\pi\mathcal O/\pi^4\mathcal O$
is cyclic, and $I=\pi^3\mathcal O$.  Nevertheless, for
$b=\pi E_{21}$,
\[
 \tau(b)b=\pi^2E_{31}\ne0
 \quad\hbox{in}\quad
 A'_{31}/(A_{32}A_{23}A_{31})
       =\pi^2\mathcal O/\pi^5\mathcal O.                    \tag{AF.4}
\]
Thus these polarization identities alone do not replace the tensor
product used in (D2.10) by an exterior square.  This example is only a
GMA calculation; it is not asserted to be a global deformation ring
or to satisfy the augmentation condition $A/I\simeq\mathcal O$.
In the paper's representation argument, Proposition~6.7 and
Lemma~6.8 (pp.~35--36) assert that $\tau$ acts identically on the
long block after the specified reduced base change.  Even such an
identity action is consistent with (AF.4); it supplies no
alternating-product identity on the integral conormal module.
There is also a separate sign issue in the displayed proof's
$M_2\times F$ case, recorded below rather than imported without
qualification.

\paragraph{The integral lower bound contains two factors, not three.}
Write, at an allowed weight $k$ of the cited paper,
\[
 A_k=L_p^\Sigma(0,\Psi_k),\qquad
 B_k=L_p^\Sigma(-1,(\Psi_+)_k\eta),\qquad
 \Psi_+=\Psi\circ\operatorname{Ver},                         \tag{AF.5}
\]
where $\eta$ is the quadratic character of the imaginary quadratic
field.  Equations~(8.3), Theorem~8.16 (p.~45), and (10.6) (p.~55)
give the actual integral surjection
\[
 T^0_k/J^0_k\longtwoheadrightarrow\mathcal O/(A_kB_k).        \tag{AF.6}
\]
Assumption (S3) and (10.7) separately give
\[
 \begin{aligned}
 a_k:=\operatorname{length}_{\mathcal O}H_{\rm adj}
       &\leq v_\pi(A_k),\\
 b_k:=\operatorname{length}_{\mathcal O}H_{\rm long}
       &\leq v_\pi(B_k),\\
 c_k:=\operatorname{length}_{\mathcal O}(T^0_k/J^0_k)
       &\geq v_\pi(A_k)+v_\pi(B_k).
 \end{aligned}                                               \tag{AF.7}
\]
Here finiteness and nonzero values are retained when writing finite
lengths.  These inequalities establish $c_k\geq a_k+b_k$.
They do not establish $c_k\geq2a_k+b_k$, the sufficient inequality
obtained from (D2.14) for two generators.  For example the numbers
$a_k=v_\pi(A_k)=c_k=1$ and $b_k=v_\pi(B_k)=0$ satisfy every
inequality in (AF.7) and fail the proposed conclusion.  This example
only checks the logical content of the bounds; it claims no arithmetic
realization.  A conjugate copy of (AF.6) does not yield a surjection
onto a quotient by $A_k^2B_k$ without an additional independence or
ideal-product statement.

Lemma~10.6 and its complete proof (pp.~57--58) do not supply that
statement.  Its second adjacent bound follows from its first by
conjugation.  Its Iwasawa specialization and local comparison give
the one-copy upper bounds, not another lower bound for the congruence
quotient.  Surplus valuation in either factor could in principle make
the stronger inequality true at some weight.  No such surplus for the
actual Tate ratios is calculated by these results.

\paragraph{The actual Tate substitution has separate scope barriers.}
Set $\chi=\chi_p$ and consider the actual characteristic-zero
triple $(\chi^e,1,\chi^{-e})$ over $G_E$, with $e>0$ even.
Matching the paper's characters
$(\Psi_k\chi^{-1},1,\Psi_k^{-c}\chi)$ gives
\[
 \Psi_k=\chi^{e+1},\qquad
 \Psi_k^{-1}\chi=\chi^{-e},\qquad
 \Psi_k^{-1-c}\chi^2=\chi^{-2e},\qquad
 (\Psi_+)_k=\chi^{2e+2}.                                    \tag{AF.8}
\]
The last equality is the transfer identity: the restriction to $G_E$
is the square of the parallel character, and at complex conjugation
both sides are one.  The long-root $c=-1$ group therefore descends
to a group over $\mathbf Q$ with coefficient
$\eta\chi^{-2e}$.  It is not another copy of the adjacent
$\chi^{-e}$ group.  For a residual Tate ratio of order five,
the powers $\omega^{-e}$ and $\omega^{-2e}$ remain distinct.
The order-five class bound and the $p\mid e$ rational strict exponent
bound do not identify these entire ordinary groups or their lengths.

Theorem~10.2 is set up in Section~10 (pp.~51--53) for infinity type
$(a+b,-b)$ with $a\geq3$ and $b\geq-1$.  The first substitution
in (AF.8) has parallel type $(e+1,e+1)$ and would require
\[
 (a,b)=(2e+2,-e-1).
                                                               \tag{AF.9}
\]
It fails $b\geq-1$.  Swapping the outer characters requires
$\Psi_k=\chi^{1-e}$ and $(a,b)=(2-2e,e-1)$, which fails
$a\geq3$.  Moreover the ordinary orders in (10.1) and (10.2)
are respectively $(\tau_1,\tau_3,1)$ and $(1,\tau_1,\tau_3)$
at the two primes over a split $p$.  Our existing source construction
does not supply these two flags.  A locally split representation may
have several flags, so no absolute impossibility is claimed.

The local conditions cannot simply be left implicit in changing the
weight.  Definition~8.22 and Example~8.23 (p.~48) impose the adjacent
inertia condition at one prime in the paper's weight cone, with the
conjugate condition at the other prime for the conjugate character.
For the parallel coefficient $\chi^{-e}$, both weights are negative,
so Definition~8.22 imposes the inertia condition at both primes.
Definition~8.24 treats the special boundary $b=-1$ and does not cover
(AF.9).  No equality of the relevant integral Selmer or conormal
modules after this change has been established.

Theorem~8.16 itself must be distinguished from the earlier unmodified
Hsieh statement.  Its proof explicitly uses Finis and He to remove
one additional finite-character ramification assumption and derives
the stated $\mu$ condition from $p\nmid\#\Delta_{\mathfrak c}$.
Thus the older Hsieh parity/ramification exclusion should not be
silently applied to this enlarged theorem.  The failed Section~10
weight and local identifications above, and the absent second
adjacent congruence bound, remain separate obstacles.  This audit
does not claim a cyclotomic factorization of the specialized Katz
value, whose interpolation, periods, and local factors would need
their own comparison.

\paragraph{A source normalization error and its ramification repair.}
The PDF at printed p.~50, equation~(9.3), visibly writes
$\#(M\otimes_{\mathcal O}\mathcal O')=[F':F]\,\#M$ for a
finite $\mathcal O$-module $M$.  With $d=[F':F]=e'f'$ the correct
formula is
\[
 \#(M\otimes_{\mathcal O}\mathcal O')=(\#M)^d,\qquad
 \operatorname{length}_{\mathcal O'}(M\otimes\mathcal O')
       =e'\operatorname{length}_{\mathcal O}M.                \tag{AF.10}
\]
These follow either from the elementary divisors or from
$\mathcal O'\simeq\mathcal O^d$ and
$\pi\mathcal O'=(\pi')^{e'}\mathcal O'$.
The error also appears in the immediately following sentence.

Here is the uniform repair of the numerical depth in Proposition~9.3,
retaining that proposition's imported Galois, ordinary, sign, and
lattice inputs.  Their production and Selmer placement of primitive
classes are not proved by this numerical repair.  Its free Hecke
rank $r$ is fixed.  In Lemma~9.2 choose
$N=rn$.  Write $q=\#k$ and $q'=q^{f'}$.  The product inequality
in (9.3), together with $\#M\geq q^{rn}$, then gives
\[
 (q')^{\sum_i m_i}\geq q^{e'f'rn}
       =(q')^{e'rn},\qquad \max_i m_i\geq e'n.               \tag{AF.11}
\]
Changing the weight back to $k$ changes its characters only at depth
$e'rn$, so the chosen congruence retains depth $e'n$ over
$\mathcal O'$.  After any further field of definition $F''/F'$
used for the representation and its stable lattice, it retains depth
\[
 e''e'n=e(F''/F)n.                                          \tag{AF.12}
\]
Use the same lattice lemma at that retained exponent.  Its primitive
extension class then has order exponent $e(F''/F)n$ with respect to
the uniformizer of $\mathcal O''$.

This really contradicts a fixed finite divisible Selmer group, even
when $F''$ depends on $n$.  Let $C$ annihilate all the finitely many
base Selmer groups in Proposition~9.3 by $\pi^C$.  Finite free
coefficient extension commutes with continuous cochain complexes,
their cohomology, and the kernels defining the inertia conditions:
after choosing an $\mathcal O$-basis of $\mathcal O''$, these are
just finite direct sums.  Also
$(F/\mathcal O)\otimes_{\mathcal O}\mathcal O''
\simeq F''/\mathcal O''$; the $c=-1$ projector commutes with
this extension since $2$ is invertible.  Thus every extended group
is killed by $\pi^C$, or equivalently by a uniformizer power
$e(F''/F)C$.  The primitive class of depth (AF.12) contradicts this
as soon as $n>C$.  The paper's finite-coefficient classes map
injectively to the divisible groups because its residual character
distinctions give $H^0(G_\Sigma,F''/\mathcal O''(\tau))=0$.
These distinctions are imposed at one element of the decomposition
group in Definition~8.19 (p.~46).
This repairs the numerical normalization of the contradiction;
no independent second adjacent factor arises from coefficient
extension.

\paragraph{A separate sign gap in the displayed algebraic argument.}
The inference in Lemma~6.8 that $\tau^2=1$ forces the implementing
matrix to be symmetric is valid for the full odd-dimensional matrix
algebra.  It is not valid for its stated third case, the outer
$M_2(F)$ block plus the middle $F$ block.  To see the precise issue,
over a field $F$ of characteristic different from two put
\[
 \mathcal B=
 \left\{\begin{pmatrix}a&0&b\\0&t&0\\c&0&d\end{pmatrix}
           :a,b,c,d,t\in F\right\},\qquad
 C=\begin{pmatrix}0&0&1\\0&1&0\\-1&0&0\end{pmatrix},
 \qquad \tau_C(X)=CX^tC^{-1}.                                \tag{AF.13}
\]
Then $\mathcal B\simeq M_2(F)\times F$, $\tau_C$ is an
anti-automorphism, and
\[
 \begin{gathered}
 C C^{-t}=\operatorname{diag}(-1,1,-1),\qquad
 \tau_C^2|_{\mathcal B}=1,\\
 \tau_C(E_{11})=E_{33},\qquad \tau_C(E_{22})=E_{22},\qquad
 \tau_C(E_{13})=-E_{13},\quad \tau_C(E_{31})=-E_{31}.
 \end{gathered}                                               \tag{AF.14}
\]
The diagonal matrix in (AF.14) lies in the centralizer of
$\mathcal B$ but is not scalar on $F^3$.  The determinant
pseudorepresentation is invariant under $\tau_C$.  Thus
determinant self-duality and outer-idempotent interchange alone
do not force the asserted positive action on the long block.
The definitions of self-duality in Section~4, including
Definitions~4.1/4.8 and Corollary~4.10, and the setup of
Section~6.1 impose those identities; no additional common
symmetric-polarization sign was found in those clauses.

This is a counterexample to the displayed algebraic inference,
not an arithmetic counterexample to Proposition~6.7 or the main
theorem.  The field-algebra model does not assert realization by
the paper's global Galois group, ordinary deformation ring, or
automorphic representations.  A prescribed symmetric ambient
polarization, an appropriate conjugate-orthogonal sign on every
irreducible block, or restriction to a full $M_3$ generic branch
would exclude it.  Supplying such a condition and proving it in
the intended application is a separate repair.  In particular,
the use of the ordinary GMA injections of Theorem~6.4 in (D2.12)--
(D2.14) remains an explicit hypothesis; this audit has not
discharged it for the actual Tate source.

\paragraph{Outcome and audit scope.}
The exact conditional source boundary in (D2.14) is unchanged.
The 2026 preprint supplies one adjacent factor and one long-root
factor in its own ordinary weight range.  Neither its polarization
identities nor its base-change cardinality argument supplies the
extra adjacent length for the actual Tate-ratio problem.  The
remaining order-five or order-at-least-seven gate stays open.
The complete cited polarization proofs, Proposition~6.7/Lemma~6.8,
Theorem~8.16, Section~8.5 definitions, Proposition~9.3,
Theorem~10.2, and Lemma~10.6 were read for this check.  Their imported
theorems are retained as inputs, not independently reproved here.
The PDF equation~(9.3) was checked visually.  The sign issue is
confined to the identified algebraic inference, with omitted
arithmetic hypotheses stated explicitly.  This is a source-scope
and normalization audit, not a Lean kernel verification or a new
automorphic vanishing result.
